\section{Assign Cookies}
\label{sec:assign-cookies}

\problemstatement{We are given two arrays. The array $g$ of length $n$ holds
the \emph{greed factor} of each child: child $i$ will be content only if it
receives a cookie of size at least $g[i]$. The array $s$ of length $m$ holds
the \emph{size} of each available cookie. Each cookie may be given to at most
one child, and each child may receive at most one cookie. We must hand out
cookies so as to \textbf{maximize the number of content children}, and return
that maximum count.}

\begin{patternbox}
The phrase \emph{``maximize the number of children you can satisfy''} together
with two independent arrays that must be \emph{matched} against each other is
the signature of a two-pointer greedy assignment problem. Whenever a problem
gives you two collections, a one-to-one matching rule of the form ``resource
$j$ can serve request $i$ if resource $j$ is large enough'', and asks for the
maximum number of satisfied requests, the standard reflex is: \emph{sort both
arrays, then match the smallest unsatisfied request with the smallest
resource capable of satisfying it.} No recursion, no DP table, and no
exhaustive matching is needed -- sorting already throws away every irrelevant
ordering and leaves only a single greedy sweep.
\end{patternbox}

\subsection*{Understanding the problem carefully}

Strip away the cookie/child story and the problem is really about two
multisets of numbers, $g$ and $s$. We are allowed to pair an element of $g$
with an element of $s$ exactly when
\[
  s[j] \;\ge\; g[i],
\]
and each element of either array can be used in at most one pair. We want the
largest possible number of valid pairs.

Notice two things immediately.

\begin{enumerate}
  \item The actual \emph{order} in which children or cookies are listed in
        the input is irrelevant to the answer -- only the multiset of greed
        factors and the multiset of cookie sizes matter. This is a strong hint
        that we are free to reorder (sort) the input without losing anything,
        which is exactly the kind of freedom a greedy algorithm wants to
        exploit.
  \item A large cookie can always satisfy any child that a smaller cookie
        could satisfy. So if we are ever tempted to ``save'' a large cookie
        for later while a smaller cookie sitting right next to it could have
        done the same job, we are never making things better, and sometimes
        we are making things strictly worse, because the large cookie might
        have been the \emph{only} cookie capable of satisfying some greedier
        child later on.
\end{enumerate}

\subsection*{Why sorting is the first move}

\begin{ideabox}[Greedy Choice]
Sort $g$ and $s$ in increasing order. Walk through the children from the
least greedy to the most greedy. For the current least-greedy unsatisfied
child, give them the \emph{smallest cookie that is still large enough} to
satisfy them. If no remaining cookie is large enough, this child -- and every
child after them, since they are all at least as greedy -- cannot be the next
one satisfied with the cookies left, so we move on without giving up a
cookie unnecessarily.
\end{ideabox}

Why does always matching the least-greedy child with the smallest sufficient
cookie never hurt? Suppose, for contradiction, that some optimal assignment
does \emph{not} give the least-greedy child $c$ the smallest sufficient
cookie $k$ available, but instead gives $c$ some larger sufficient cookie
$k'$, while $k$ is given to a different child $c'$ (or not used at all). Since
$c$ is the least greedy child being considered, and $k$ is already known to
satisfy $c$ (it is sufficient), and $k$ is the smallest such cookie, swapping
$k$ and $k'$ between $c$ and $c'$ keeps $c$ satisfied (since $k$ is still
sufficient for $c$ by definition) and does not break $c'$'s satisfaction
either, because we only handed $c'$ a \emph{larger} cookie $k'$ than before --
and a larger cookie satisfies anyone a smaller one would. So the swap never
decreases the number of satisfied children. This is a standard
\emph{exchange argument}: any optimal solution can be rearranged, one swap at
a time, into the greedy solution without ever lowering the count of satisfied
children. That is precisely what makes the greedy choice safe.

\subsection*{Why two pointers are enough}

Once both arrays are sorted, we never need to look backwards. We keep one
pointer $i$ over children (starting at the least greedy) and one pointer $j$
over cookies (starting at the smallest cookie).

\begin{itemize}
  \item If the current cookie $s[j]$ is large enough for the current child
        $g[i]$, this is the cheapest cookie that can possibly satisfy this
        child (every cookie before $j$ was too small, by sorted order), so we
        commit: give it to the child, and advance both $i$ and $j$.
  \item If $s[j]$ is too small for $g[i]$, this cookie is too small for
        \emph{every} remaining child too (they are all at least as greedy as
        $g[i]$), so it can never be used -- discard it by advancing $j$ alone.
\end{itemize}

Each pointer only ever moves forward, and each comparison either satisfies a
child or permanently discards a cookie. There is no benefit to revisiting a
discarded cookie or an already-satisfied child, so a single linear sweep after
sorting is provably sufficient -- no recursion or dynamic programming is
needed because there is no overlapping subproblem to memoize; every decision
is final the moment it is made.

\subsection*{Algorithm}

\begin{enumerate}
  \item Sort $g$ in increasing order; sort $s$ in increasing order.
  \item Set $i \gets 0$ (pointer into $g$), $j \gets 0$ (pointer into $s$),
        $\textit{count} \gets 0$.
  \item While $i < n$ and $j < m$:
  \begin{enumerate}
    \item If $s[j] \ge g[i]$: this cookie satisfies this child. Increment
          \textit{count}, and advance both $i \gets i+1$ and $j \gets j+1$.
    \item Otherwise: this cookie is too small for every remaining child.
          Advance $j \gets j+1$ only.
  \end{enumerate}
  \item Return \textit{count}.
\end{enumerate}

\begin{algorithm}[H]
\caption{Assign Cookies}
\begin{algorithmic}[1]
\Procedure{FindContentChildren}{$g, s$}
  \State sort $g$ in increasing order
  \State sort $s$ in increasing order
  \State $i \gets 0$, $j \gets 0$, $\textit{count} \gets 0$
  \While{$i < \text{length}(g)$ \textbf{and} $j < \text{length}(s)$}
    \If{$s[j] \ge g[i]$}
      \State $\textit{count} \gets \textit{count} + 1$
      \State $i \gets i + 1$
      \State $j \gets j + 1$
    \Else
      \State $j \gets j + 1$
    \EndIf
  \EndWhile
  \State \Return \textit{count}
\EndProcedure
\end{algorithmic}
\end{algorithm}

\subsection*{Dry run}

\dryrun{Take $g = [1, 2, 3]$ and $s = [1, 1]$.}

Both arrays are already sorted. We initialize $i = 0$, $j = 0$,
$\textit{count} = 0$.

\begin{itemize}
  \item $i=0, j=0$: compare $s[0]=1$ with $g[0]=1$. Since $1 \ge 1$, the
        child is satisfied. $\textit{count} = 1$, $i \to 1$, $j \to 1$.
  \item $i=1, j=1$: compare $s[1]=1$ with $g[1]=2$. Since $1 < 2$, this
        cookie is too small. $j \to 2$.
  \item $j = 2$ equals the length of $s$, so the loop ends.
\end{itemize}

The final answer is $\textit{count} = 1$. Indeed, with cookies of size
$[1,1]$ we can satisfy only the child with the smallest greed factor $1$; the
children needing $2$ and $3$ can never be satisfied by a cookie of size $1$.

Now take a second example: $g = [1, 2]$ and $s = [1, 2, 3]$.

\begin{itemize}
  \item $i=0, j=0$: $s[0]=1 \ge g[0]=1$. Satisfied. $\textit{count}=1$,
        $i \to 1$, $j \to 1$.
  \item $i=1, j=1$: $s[1]=2 \ge g[1]=2$. Satisfied. $\textit{count}=2$,
        $i \to 2$, $j \to 2$.
  \item $i = 2$ equals the length of $g$, so the loop ends.
\end{itemize}

The final answer is $\textit{count} = 2$: every child is satisfied, and the
leftover cookie of size $3$ is simply not needed.

\subsection*{C++ implementation}

\begin{lstlisting}
#include <bits/stdc++.h>
using namespace std;

int findContentChildren(vector<int>& g, vector<int>& s) {
    sort(g.begin(), g.end());
    sort(s.begin(), s.end());

    int i = 0, j = 0, count = 0;
    int n = g.size(), m = s.size();

    while (i < n && j < m) {
        if (s[j] >= g[i]) {
            count++;
            i++;
            j++;
        } else {
            j++;
        }
    }
    return count;
}

int main() {
    vector<int> g = {1, 2, 3};
    vector<int> s = {1, 1};

    cout << "Number of content children: "
         << findContentChildren(g, s) << endl;
    return 0;
}
\end{lstlisting}

\begin{complexitybox}
\textbf{Time Complexity.} Sorting $g$ takes $O(n \log n)$ and sorting $s$
takes $O(m \log m)$. The subsequent two-pointer sweep visits each index of
$g$ and $s$ at most once, so it costs $O(n + m)$. The sorting step dominates,
giving an overall time complexity of
\[
  O(n \log n + m \log m).
\]
\textbf{Space Complexity.} Sorting can be done in place, and we use only the
two pointers $i, j$ and the counter as extra variables, so the auxiliary
space complexity is $O(1)$ (ignoring the in-place sort's recursion stack,
which is $O(\log n)$ in the worst case for typical \texttt{std::sort}
implementations).
\end{complexitybox}

\begin{takeawaybox}
When a problem asks you to \emph{match} two sorted resources against each
other to maximize the count of satisfied requirements, sort both sides and
sweep with two pointers, always trying to satisfy the smallest unsatisfied
requirement with the cheapest resource capable of doing so. This pattern --
\emph{sort, then greedily match smallest-feasible-to-smallest-need} -- recurs
throughout greedy algorithms and is worth recognizing on sight.
\end{takeawaybox}
